\section{Conclusiones}

En esta práctica aplicamos PSO al problema de encontrar la conformación de mínima energía en el modelo HP bidimensional. El modelo reduce una proteína a una cadena de aminoácidos hidrofóbicos y polares sobre una cuadrícula, y aun así el problema es NP-completo y su espacio de búsqueda crece como $2.638^n$. La dificultad principal fue que PSO trabaja con vectores continuos y una conformación es un objeto discreto que no puede cruzarse consigo mismo.

La representación que elegimos resolvió ese problema de forma razonable. Usar direcciones relativas redujo el espacio de búsqueda y eliminó las rotaciones equivalentes. Guardar tres preferencias por paso, en lugar de un solo número redondeado, evitó imponer un orden artificial entre girar a la izquierda, seguir de frente y girar a la derecha, y le dio a la actualización de PSO una interpretación directa: las partículas suben la preferencia por las decisiones de sus mejores conformaciones. El decodificador con reparación fue indispensable. En cadenas de solo 14 aminoácidos, menos de la mitad de las descripciones al azar son conformaciones válidas, y la fracción baja con la longitud; con la reparación, todas las ejecuciones terminaron en conformaciones completas y válidas.

En las cadenas cortas el algoritmo funcionó bien. Encontró la energía mínima conocida en 20 de 30 ejecuciones en S1-20, y al menos una vez en S2-24 y en S4-36. En las cadenas de 48 y 64 aminoácidos no llegó al óptimo: el error relativo de la media fue de 18.6\,\% y 29.2\,\%, y la mejor ejecución quedó a dos y nueve contactos de la mejor energía conocida. Las conformaciones encontradas sí tienen la estructura esperada, con un núcleo de H rodeado de P, pero les faltan contactos.

Las medidas de diversidad explicaron esa diferencia mejor que el \textit{gbest} solo. En las cadenas cortas las tres medidas cayeron casi a cero antes de la iteración 750, y en S1-20 el \textit{gbest} dejó de mejorar, en promedio, en la iteración 59: el enjambre colapsó sobre una conformación y el resto de las iteraciones no aportó nada. En las cadenas largas la diversidad seguía bajando al final y la última mejora ocurrió, en promedio, después de la iteración 550. Ahí el enjambre se quedó sin iteraciones antes de terminar de concentrarse, así que el siguiente paso más directo sería darle más tiempo a esas cadenas. La diversidad cognitiva fue la que más tardó en bajar en las cadenas largas, lo que indica que las memorias de las partículas conservan regiones distintas incluso cuando sus posiciones actuales ya se agruparon.

La estrategia de coeficientes con curva coseno y recalentamiento fue la mejor de las cuatro que probamos, con diferencias significativas contra las curvas lineal y exponencial en las dos cadenas, y contra el coseno sin recalentamiento en S4-36. El orden de las estrategias coincidió con cuánto tiempo mantienen la exploración: la exponencial, que la reduce primero, fue la peor. Este resultado, junto con el error que encontramos en la primera versión del recalentamiento (que dejaba $w$ casi fijo en 0.88 y aun así funcionaba bien), nos lleva a la misma conclusión: en un paisaje lleno de mesetas como el del modelo HP, conviene explorar durante mucho tiempo. También aprendimos que conviene graficar también los coeficientes que realmente se usan, porque el error de la primera versión solo se veía en esa gráfica.

Quedan varias mejoras por probar. La más prometedora es combinar PSO con una búsqueda local sobre el \textit{gbest}, por ejemplo con \textit{pull moves}~\cite{lesh2003complete}, que es lo que hacen los métodos que alcanzan las mejores energías conocidas en cadenas largas. También valdría la pena probar una topología de vecindario en anillo, que tiende a conservar más diversidad que la topología global, y repetir los experimentos con más iteraciones en las cadenas de 48 aminoácidos o más.
